\documentclass{article}
\usepackage{graphicx} % Required for inserting images
\usepackage{amsmath}
\usepackage{amssymb}
\usepackage{amsfonts}

\title{Notes on "An Introduction to Manifolds" by Tu}
\author{Kevin Sony}
\date{Spring 2026}

\newcommand{\norm}[1]{\left\lVert#1\right\rVert}

\newcommand{\betabar}{\bar{\beta}}
\newcommand{\kbar}{\bar{k}}
\newcommand{\R}{\mathbb{R}}
\newcommand{\Integers}{\mathbb{Z}}
\newcommand{\Cinfty}{C^{\infty}}
\newcommand{\comment}[1]{}


\begin{document}
	
	\maketitle
	
	\section{Smooth Functions on Euclidean Space}
	In an introductory calculus course, we had worked with $\R^n$. $\R^n$ is a manifold, but it also is equipped with a natural coordinate system, unlike many other manifolds. It turns out that a lot of stuff from calculus that was done was only possible because there existed a corresponding global coordinate system. For example, it is impossible to integrate functions on a manifold in general, since integration requires a coordinate system. The objects that can be integrated are differential forms.
	
	The goal is to recast calculus in $\R^n$ in a way independent of coordinates to generalize to manifolds.
	
	\subsection{Smooth vs Analytic Functions}
	Write coordinates on $\R^n$ as $x^1, \cdots, x^n$ and let $p = (p^1, \cdots, p^n)$ be a point in an open set of $U \subset \R^n$. Index using indices as superscripts (there is reason for this).
	
	\subsubsection*{Definition}
	Let $k \in \Integers^+_0$. A real-valued function $f: U \to \R^n$ is said to be $C^k$ at $p \in U$ if its partial derivatives
	\begin{gather*}
		\dfrac{\partial^j f}{\partial x^{i_1} \cdots \partial x^{i_j}}
	\end{gather*}
	of all orders $j \leq k$ exists and are continuous at $p$. The function $f: U \to \R^n$ is $C^{\infty}$ at $p \in U$ if it is $C^k$ for all $k \geq 0$. A vector-valued function $f: U \to \R^m$ is said to be $C^k$ if all of its component functions $f^1, \cdots, f^m$ are $C^k$ at $p$. We say that $f: U \to \R^m$ is $C^k$ on $U$ if it is $C^k$ for every point $p \in U$. The definition for $C^{\infty}$ for vector-valued functions are defined in a similar manner. $C^{\infty}$ and smooth are equivalent.
	
	A neighborhood of a point in $\R^n$ is an open set containing the point. The function $f$ is real analytic at $p$ if in some neighborhood of $p$ it is equal to its Taylor series at $p$, i.e.
	\begin{gather*}
		f(x) = f(p) + \sum_i \dfrac{\partial f}{\partial x^i}(p) (x^i - p^i) + \dfrac{1}{2!} \sum_{i, j} \dfrac{\partial^2 f}{\partial x^i \partial x^j}(p) (x^i - p^i) (x^j - p^j) \\
		+ \cdots + \dfrac{1}{k!} \sum_{i_1, \cdots, i_k} \dfrac{\partial^k f}{\partial x^{i_1} \cdots \partial x^{i_j}}(p) (x^{i_1} - p^{i_1}) \cdots (x^{i_k} - p^{i_k}) + \cdots
	\end{gather*} 
	
	All real-analytic functions are $C^{\infty}$, but not all $C^{\infty}$ functions are real-analytic. Consider the function
	\begin{gather*}
		f(x) = \begin{cases}
			e^{-1/x}, \quad &x > 0 \\
			0, \quad &x \leq 0
		\end{cases}
	\end{gather*}
	It is relatively easy to show via induction that $f(x)$ is $C^{\infty}$. However, note that at $x = 0$, $f(0) = 0$ and $f^{(k)}(0) = 0$ and so its Taylor series expansion should be $0$. However, for $x > 0$, $f(x) = e^{-1/x} \neq 0$ and so there does not exist a neighborhood around $x = 0$ such that $f(x)$ is identical to its Taylor series.
	
	\subsection{Taylor's Theorem with Remainders}
	While a $C^{\infty}$ is not necessarily equal to its Taylor series, there is a Taylor's theorem with remainder that is good enough. In the theorem below, the Taylor series will only consist of the constant term $f(p)$
	
	\subsubsection*{Taylor's theorem with remainders}
	Let $f$ be $C^{\infty}$	function on an open convex subset $U \subset \R^n$. Let $p = (p^1, \cdots, p^n) \in U$. Then there exists functions $g_1(x), \cdots, g_n(x) \in C^{\infty}(U)$ such that
	\begin{gather*}
		f(x) = f(p) + \sum_{i=1}^{n} (x^i - p^i) g_i(x)
	\end{gather*}
	where
	\begin{gather*}
		g_i(p) = \dfrac{\partial f}{\partial x^i}(p)
	\end{gather*}

	\noindent \textit{Proof}: For any point $x \in U$, there exists a line segment entirely in $U$ which connects $x$ to $p$, given by $p + t(x - p)$, for $0 \leq t \leq 1$. Thus, $f(p + t(x - p))$ exists and is $C^{\infty}$ for all points along the line segment. Taking the derivative of this function with respect to $t$, we get
	\begin{gather*}
		\dfrac{d}{dt} f(p + t(x - p)) = \sum_{i=1}^{n} (x^i - p^i) \dfrac{\partial f}{\partial x^i} (p + t(x - p))
	\end{gather*}
	Integrating with respect to $t$ from 0 to 1, we get
	\begin{gather*}
		\int_{0}^{1} \dfrac{d}{dt} f(p + t(x - p)) dt = \int_{0}^{1} \sum_{i=1}^{n} (x^i - p^i) \dfrac{\partial f}{\partial x^i} (p + t(x - p)) dt
	\end{gather*}
	The LHS evaluates to
	\begin{gather*}
		\int_{0}^{1} \dfrac{d}{dt} f(p + t(x - p)) dt = f(p + t(x - p)) \Big|_0^1
	\end{gather*}
	and the RHS evaluates to
	\begin{gather*}
		\int_{0}^{1} \sum_{i=1}^{n} (x^i - p^i) \dfrac{\partial f}{\partial x^i} (p + t(x - p)) dt = \sum_{i=1}^{n} (x^i - p^i) \int_{0}^{1} \dfrac{\partial f}{\partial x^i} (p + t(x - p)) dt
	\end{gather*}
	If we define $g_i(x)$ to be
	\begin{gather*}
		g_i(x) = \int_{0}^{1} \dfrac{\partial f}{\partial x^i} (p + t(x - p)) dt
	\end{gather*}
	we see that $g_i(x)$ is $C^{\infty}$ as $f$ is also $C^{\infty}$. Combining the results, we have that
	\begin{gather*}
		f(x) - f(p) = \sum_{i=1}^{n} (x^i - p^i) g_i(x)
	\end{gather*}
	Now finally, notice that when $x = p$, $p + t(x - p) = p + t(p - p) = p$ and so
	\begin{gather*}
		g_i(p) = \int_{0}^{1} \dfrac{\partial f}{\partial x^i} (p) dt = \dfrac{\partial f}{\partial x^i} (p) \int_{0}^{1} dt = \dfrac{\partial f}{\partial x^i} (p)
	\end{gather*}
	
	\subsubsection*{Exercise 1.2}
	Let $f(x)$ be
	\begin{gather*}
		f(x) = \begin{cases}
			e^{-1/x}, \quad &x > 0 \\
			0, \quad &x \leq 0
		\end{cases}
	\end{gather*}
	Show by induction that for $x > 0$ and $k \geq 0$, the $k$th derivative $f^{(k)}(x)$ is of the form $p_{2k}(1/x)e^{-1/x}$ for some polynomial $p_{2k}(y)$ of degree $2k$ in $y$. Prove that $f$ is $C^{\infty}$ on $\R$ and that $f^{(k)}(0) = 0$ for all $k \geq 0$
	
	Clearly, for $x > 0$, $f(x) = e^{-1/x}$. In the base case $n = 0$, $f(x) = e^{-1/x}$ and so $p_0 = 1$ is a polynomial of degree 0. Now suppose that for some $n = k$, $f^{(k)}(x) = p_{2k}(1/x)e^{-1/x}$ for some polynomial $p_{2k}(y)$ of degree $2k$ in $y$. Then
	\begin{gather*}
		f^{(k + 1)}(x) = \left( p_{2k}'(1/x) \cdot -\dfrac{1}{x^2} e^{-1/x} \right) + \left( p_{2k}(1/x)e^{-1/x} \cdot \dfrac{1}{x^2} \right) \\
		= e^{-1/x} \left( 1 / x \right)^2 \left( p_{2k}(1/x) - p_{2k}'(1/x) \right)
	\end{gather*}
	Letting $y = 1/x$, we see that
	\begin{gather*}
		e^{-1/x} \left( 1 / x \right)^2 \left( p_{2k}(1/x) - p_{2k}'(1/x) \right) = e^{-1/x} y^2 \left( p_{2k}(y) - p_{2k}'(y) \right) \\
		= e^{-1/x} p_{2k + 2}(y) = e^{-1/x} p_{2(k + 1)}(1/x)
	\end{gather*}
	Self-evidently, when $x < 0$, $f(x) = 0$ is $C^{\infty}$ and likewise, when $x > 0$, $f(x) = e^{-1/x}$ is $C^{\infty}$. The only contention arises when dealing with $x = 0$, as one must show all the derivatives of $f(x)$ is continuous there.
	
	To show that $f(x)$ is $C^{\infty}$ at 0, we need to show that for $k \geq 0$
	\begin{gather*}
		\lim_{x \to 0} f^{(k)}(x) = 0
	\end{gather*}
	Clearly, since $f(x) = 0$ for $x < 0$, we have that
	\begin{gather*}
		\lim_{x \to 0^-} f^{(k)}(x) = 0
	\end{gather*}
	Since $f(x) = e^{-1/x}$, $f^{(k)}(x) = p_{2k}(1/x)e^{-1/x}$ for some polynomial $p_{2k}(y)$ of degree $y$.
	\begin{gather*}
		\lim_{x \to 0^+} f^{(k)}(x) = \lim_{y \to \infty} p_{2k}(y)e^{-y} = 0
	\end{gather*}
	Therefore, $f(x)$ is $C^{\infty}$ at 0.

	\subsubsection*{Exercise 1.6}
	Prove that if $f: \R^2 \to \R$ is $C^{\infty}$, then there exists functions $g_{11}, g_{12}, g_{22}$ such that
	\begin{gather*}
		f(x, y) = f(0, 0) + \dfrac{\partial f}{\partial x}(0, 0)x + \dfrac{\partial f}{\partial y}(0, 0)y + x^2 g_{11}(x, y) \\
		+ xy g_{12}(x, y) + y^2 g_{22}(x, y)
	\end{gather*}
	
	Since we know that $f$ is $C^{\infty}$, we know from Taylor's theorem with Remainder that there exists $C^{\infty}$ functions $g_1, g_2$ such that
	\begin{gather*}
		f(x, y) = f(0, 0) + (x - 0) g_1(x, y) + (y - 0) g_2(x, y) \\
		= f(0, 0) + x g_1(x, y) + y g_2(x, y)
	\end{gather*}
	and
	\begin{gather*}
		g_1(0, 0) = \dfrac{\partial f}{\partial x} (0, 0) \\
		g_2(0, 0) = \dfrac{\partial f}{\partial y} (0, 0)
	\end{gather*}
	Applying Taylor's theorem with remainder to the function $g_1(x, y)$, we see that there exists functions $h_{11}, h_{12}$ in $C^{\infty}$ such that
	\begin{gather*}
		g_1(x, y) = g_1(0, 0) + (x - 0) h_{11}(x, y) + (y - 0) h_{12}(x, y) \\
		= \dfrac{\partial f}{\partial x} (0, 0) + x h_{11}(x, y) + y h_{12}(x, y)
	\end{gather*}
	Similarly, applying it to $g_2(x, y)$, we see that there exists functions $h_{21}, h_{22}$ in $C^{\infty}$ such that
	\begin{gather*}
		g_2(x, y) = g_2(0, 0) + (x - 0) h_{21}(x, y) + (y - 0) h_{22}(x, y) \\
		= \dfrac{\partial f}{\partial y} (0, 0) + x h_{21}(x, y) + y h_{22}(x, y)
	\end{gather*}
	Plugging these results in, we have that
	\begin{gather*}
		f(x, y) = f(0, 0) + x \left( \dfrac{\partial f}{\partial x} (0, 0) + x h_{11}(x, y) + y h_{12}(x, y) \right) \\
		+ y \left( \dfrac{\partial f}{\partial y} (0, 0) + x h_{21}(x, y) + y h_{22}(x, y) \right) \\
		f(x, y) = f(0, 0) + \dfrac{\partial f}{\partial x} (0, 0) x + \dfrac{\partial f}{\partial y} (0, 0) y + x^2 h_{11}(x, y)
		\\ + xy \left( h_{12}(x, y) + h_{21}(x, y) \right) + y^2 h_{22}(x, y)
	\end{gather*}
	Now let $g_{11}(x, y) = h_{11}(x, y)$, $g_{12}(x, y) = h_{12}(x, y) + h_{21}(x, y)$, and $g_{22}(x, y) = h_{22}(x, y)$. This completes the proof.

	
	\subsubsection*{Exercise 1.7}
	Let $f: \R^2 \to \R$ be a $C^{\infty}$ function with
	\begin{gather*}
		f(0, 0) = \dfrac{\partial f}{\partial x}(0, 0) = \dfrac{\partial f}{\partial y}(0, 0) = 0
	\end{gather*}
	Define a function $g: \R^2 \to \R$ by
	\begin{gather*}
		g(t, u) = \begin{cases}
			\dfrac{f(t, tu)}{t}, \quad &t \neq 0 \\
			0, \quad &t = 0
		\end{cases}
	\end{gather*}
	Prove that $g(t, u)$ is $C^{\infty}$ for $(t, u) \in \R^2$
	
	Recall from the earlier exercise 1.6 that since $f$ is $C^{\infty}$, then there exists functions $g_{11}, g_{12}, g_{22}$ such that
	\begin{gather*}
		f(x, y) = f(0, 0) + \dfrac{\partial f}{\partial x}(0, 0)x + \dfrac{\partial f}{\partial y}(0, 0)y + x^2 g_{11}(x, y) \\
		+ xy g_{12}(x, y) + y^2 g_{22}(x, y)
	\end{gather*}
	Plugging in the earlier conditions, we have that
	\begin{gather*}
		f(x, y) = x^2 g_{11}(x, y) + xy g_{12}(x, y) + y^2 g_{22}(x, y)
	\end{gather*}
	For $x = t, y = tu$, we have that
	\begin{gather*}
		f(t, tu) = t^2 g_{11}(t, tu) + t^2u g_{12}(t, tu) + (tu)^2 g_{22}(t, tu)
	\end{gather*}	
	and so finally, we have that
	\begin{gather*}
		\dfrac{f(t, tu)}{t} = t g_{11}(t, tu) + tu g_{12}(t, tu) + tu^2 g_{22}(t, tu) \\
		= t \left( g_{11}(t, tu) + u g_{12}(t, tu) + u^2 g_{22}(t, tu) \right)
	\end{gather*}
	Recall the earlier definition of $g(t, u)$. The derivative of this function with respect to $t$ is
	\begin{gather*}
		\dfrac{\partial g}{\partial t}(t, u) = \left( g_{11}(t, tu) + u g_{12}(t, tu) + u^2 g_{22}(t, tu) \right) \\
		+ t \left(
		\dfrac{\partial g_{11}}{\partial x} (t, tu) \dfrac{\partial x}{\partial t} + \dfrac{\partial g_{11}}{\partial y} (t, tu) \dfrac{\partial y}{\partial t}
		+ u \left( \dfrac{\partial g_{12}}{\partial x} (t, tu) \dfrac{\partial x}{\partial t} + \dfrac{\partial g_{12}}{\partial y} (t, tu) \dfrac{\partial y}{\partial t} \right) \right) \\
		+ t u^2 \left( \dfrac{\partial g_{22}}{\partial x} (t, tu) \dfrac{\partial x}{\partial t} + \dfrac{\partial g_{22}}{\partial y} (t, tu) \dfrac{\partial y}{\partial t} \right) \\
		= \left( g_{11}(t, tu) + u g_{12}(t, tu) + u^2 g_{22}(t, tu) \right) \\
		+ t \left(
		\dfrac{\partial g_{11}}{\partial x} (t, tu) + u \dfrac{\partial g_{11}}{\partial y} (t, tu)
		+ u \left( \dfrac{\partial g_{12}}{\partial x} (t, tu) + u \dfrac{\partial g_{12}}{\partial y} (t, tu) \right) \right) \\
		+ tu^2 \left( \dfrac{\partial g_{22}}{\partial x} (t, tu) + u \dfrac{\partial g_{22}}{\partial y} (t, tu) \right)
	\end{gather*}
	Similarly, the derivative of $g$ with respect to $u$ is
	\begin{gather*}
		\dfrac{\partial g}{\partial u}(t, u) = t \left(
		\dfrac{\partial g_{11}}{\partial x} (t, tu) \dfrac{\partial x}{\partial u} + \dfrac{\partial g_{11}}{\partial y} (t, tu) \dfrac{\partial y}{\partial u}
		+ g_{12}(t, tu) + u \left( \dfrac{\partial g_{12}}{\partial x} (t, tu) \dfrac{\partial x}{\partial u} + \dfrac{\partial g_{12}}{\partial y} (t, tu) \dfrac{\partial y}{\partial u} \right) \right) \\
		+ t \left( 2u g_{22}(t, tu) + u^2 \left( \dfrac{\partial g_{22}}{\partial x} (t, tu) \dfrac{\partial x}{\partial u} + \dfrac{\partial g_{22}}{\partial y} (t, tu) \dfrac{\partial y}{\partial u} \right) \right) \\
		= t \left(
		t \dfrac{\partial g_{11}}{\partial y} (t, tu)
		+ g_{12}(t, tu) + u \left( t \dfrac{\partial g_{12}}{\partial y} (t, tu) \right) + 2u g_{22}(t, tu) + u^2 \left( t \dfrac{\partial g_{22}}{\partial y} (t, tu) \right) \right)
	\end{gather*}
	Notice where $t \neq 0$
	\begin{gather*}
		g(t, u) = \dfrac{f(t, tu)}{t} = t \left( g_{11}(t, tu) + u g_{12}(t, tu) + u^2 g_{22}(t, tu) \right)
	\end{gather*}
	is $C^{\infty}$ for all points $(t, u)$ where $t \neq 0$ since all the component parts are also $C^{\infty}$. It is also continuous at $t = 0$ since
	\begin{gather*}
		\lim_{t \to 0} t \left( g_{11}(t, tu) + u g_{12}(t, tu) + u^2 g_{22}(t, tu) \right) = 0
	\end{gather*}
	since $g_{11}, g_{12}, g_{22}$ are $C^{\infty}$ and so must converge to a definite value. Finally, notice that the partial derivatives
	\begin{gather*}
		\nabla g(t, u) = \left( \dfrac{\partial g}{\partial t}(t, u), \dfrac{\partial g}{\partial u}(t, u) \right)
	\end{gather*}
	are composed of functions that at the limit is $C^{\infty}$ at $t = 0$ and so the partial derivatives are $C^{\infty}$ at $t = 0$.
	
	
	\section*{Tangent vectors in $\R^n$ as Derivations}
	We typically express a vector $v$ at a point $p \in \R^n$ as a column of numbers
	\begin{gather*}
		v = \begin{bmatrix}
			v^1 \\
			\vdots \\
			v^n
		\end{bmatrix}
	\end{gather*}
	or as a row of numbers bounded by angular brackets
	\begin{gather*}
		v = \langle v^1, \cdots, v^n \rangle
	\end{gather*}
	Recall that a secant plane to a surface in $\R^3$ is given by three non-collinear points on a surface, and as the three points approach a point $p$ on the surface, the plane at the limit position of the secant plane is called the tangent plane at $p$. Intuitively, we can understand this as the plane simply "touching" the surface at this point. A vector at $p$ is tangent to the surface if it lies in the tangent plane at $p$.
	
	
	\subsection*{Directional Derivative}
	In $\R^n$, we can visualize the tangent space $T_p(\R^n)$ at $p$ as the vector space of all vectors emanating from $p$. Since we can express the vector space with a list of $n$ numbers, it is identical to $\R^n$. To distinguish between points and vectors then, we let a point $p \in \R^n$ be written as $p = \left( p^1, \cdots, p^n \right)$ and a vector $v \in T_p(\R^n)$ be written as
	\begin{gather*}
		v = \begin{bmatrix}
			v^1 \\
			\vdots \\
			v^n
		\end{bmatrix} \quad \text{or} \quad 
		v = \langle v^1, \cdots, v^n \rangle
	\end{gather*}
	Typically, the standard basis for $\R^n$ or $T_p(\R^n)$ is denoted by $e_1, \cdots, e_n$. Then, we can express $v \in T_p(\R^n)$ as $v = \sum v^i e_i$ for some $v^i \in \R$. Elements of $T_p(\R^n)$ are called tangent vectors at $p$ in $\R^n$.
	
	The line through a point $p = \left( p^1, \cdots, p^n \right)$ with a direction $v = \langle v^1, \cdots, v^n \rangle$ has parametrization
	\begin{gather*}
		c(t) = \left( p^1 + t v^1, \cdots, p^n + t v^n \right)
	\end{gather*}
	If $f$ is $C^{\infty}$ at $p$ and $v$ is a tangent vector at $p$, then the directional derivative of $f$ in the direction of $v$ at $p$ is defined to be
	\begin{gather*}
		D_v f = \lim_{t \to 0} \dfrac{f(c(t)) - f(p)}{t} = \dfrac{d}{dt} \Bigg|_{t = 0} f(c(t))
	\end{gather*}
	By the chain rule,
	\begin{gather*}
		D_v f = \sum_{i = 1}^n \dfrac{dc^i}{dt}(0) \dfrac{\partial f}{\partial x^i}(p)
	\end{gather*}
	We know that $c'(t) = v$ and so
	\begin{gather*}
		D_v f = \sum_{i = 1}^n v^i \dfrac{\partial f}{\partial x^i}(p) = v \cdot \nabla f(p)
	\end{gather*}
	\textit{The reason the notation $D_v f$ is used, instead of $D_v f(p)$, is that the point $p$ is already encoded in the vector $v$, since $v \in T_p \R^n$. $v$ is anchored at $p$. Thus, the directional derivative is inherently taken at $p$. Moreover, $D_v f$ is understood as a number, the directional derivative at $p$, and not a function of $p$.}
	
	We write
	\begin{gather*}
		D_v = \sum_{i = 1}^n v^i \dfrac{\partial}{\partial x^i} \Bigg|_p
	\end{gather*}
	to be the map that sends a function $f$ to the number $D_v f$. The association $v \mapsto D_v$ of the directional derivative $D_v$ to the tangent vector $v$ gives a way to characterize tangent vectors as operators on functions.
	
	
	\subsection*{Germs of Functions}
	
	\subsubsection*{Equivalence Relation}
	A relation on a set $S$ is a subset $R \subset S \times S$. $R$ is an equivalence relation if it satisfies for all $x, y, z \in S$
	\begin{enumerate}
		\item (Reflexivity) $x \sim x$
		\item (Symmetry) $x \sim y$ if and only if $y \sim x$
		\item (Transitivity) $x \sim y$ and $y \sim z$ implies $x \sim z$
	\end{enumerate}
	
	If two functions agree on some neighborhood of a point $p$, they would have the same directional derivatives at $p$, suggesting that an equivalence relation can be defined on the $\Cinfty$ functions defined on some neighborhood of $p$.
	
	\subsubsection*{Germs of Functions}
	Consider the set of all pairs $(f, U)$, where $U$ is a neighborhood of $p$ and $f: U \to \R$ is $\Cinfty$. We say that two pairs $(f, U)$ and $(g, V)$ are equivalent to each other if there exists an open set $W \subset U \cap V$ containing $p$ such that $f = g$ when restricted to $W$. The equivalence class of $(f, U)$ is called the germ of $f$ at $p$. We write $\Cinfty_p(\R^n)$ (or $\Cinfty_p$ if there is no possibility of confusion) for the set of all germs of $\Cinfty$ functions on $\R^n$ at $p$.
	
	Consider the function
	\begin{gather*}
		f(x) = \dfrac{1}{1 - x}
	\end{gather*}
	over the domain $\R - \{ 1 \}$ and
	\begin{gather*}
		g(x) = 1 + x + x^2 + x^3 + \cdots
	\end{gather*}
	with the domain $(-1, 1)$. $f$ and $g$ have the same germ at any point $p \in (-1, 1)$.
	
	\subsubsection*{Algebra}
	An algebra over a field $K$ is a vector space $A$ over $K$ with the multiplication map
	\begin{gather*}
		\mu: A \times A \to A
	\end{gather*}
	denoted typically by $\mu(a, b) = a \cdot b$, such that for all $a, b, c \in A$ and $r \in K$
	\begin{enumerate}
		\item (Associativity) $(a \cdot b) \cdot c = a \cdot (b \cdot c)$
		\item (Distribution) $a \cdot (b + c) = a \cdot b + a \cdot c$ and $(a + b) \cdot c = a \cdot c + a \cdot b$
		\item (Homegeneuity) $r(a \cdot b) = (ra) \cdot b = a \cdot (rb)$
	\end{enumerate}
	This is equivalent to stating that an algebra over a field $K$ is a ring $A$ that is also a vector space over $K$ such that the ring multiplication satisfies the homogeneity condition. 

	\subsubsection*{Linear Operator}
	A map $L: V \to W$ between vector spaces over a field $K$ is called a linear operator if for any $r \in K$ and $u, v \in V$
	\begin{enumerate}
		\item $L(u + v) = L(u) + L(v)$
		\item $L(rv) = rL(v)$
	\end{enumerate}
	
	If $A$ and $A'$ are algebras over a field $K$, then an algebra homomorphism is a linear map $L: A \to A'$ that preserves the algebra multiplication, i.e. $L(a \cdot b) = L(a) \cdot L(b)$ for all $a, b \in A$.
	
	The addition and multiplication of functions induces the corresponding operations on $\Cinfty_p$, making it into an algebra over $\R$.
	
	\subsection*{Derivation at a Point}
	For each tangent vector $v$ at a point $p \in \R^n$, the directional derivative at $p$ gives a map of real vector spaces
	\begin{gather*}
		D_v: \Cinfty_p \to \R
	\end{gather*}
	Consider two functions $f, g \in \Cinfty_p$ and a scalar $r \in \R$. Recall that since
	\begin{gather*}
		D_v = \sum_{i = 1}^n v^i \dfrac{\partial}{\partial x^i} \Bigg|_p
	\end{gather*}
	we have that
	\begin{gather*}
		D_v (f + g) = \sum_{i = 1}^n v^i \dfrac{\partial (f + g)}{\partial x^i}(p) = \sum_{i = 1}^n v^i \dfrac{\partial f}{\partial x^i}(p) + \sum_{i = 1}^n v^i \dfrac{\partial g}{\partial x^i}(p) \\
		= D_v f + D_v g
	\end{gather*}
	Similarly,
	\begin{gather*}
		D_v (rf) = \sum_{i = 1}^n v^i \dfrac{\partial (rf)}{\partial x^i}(p) = r \sum_{i = 1}^n v^i \dfrac{\partial f}{\partial x^i}(p) = r D_v f
	\end{gather*}
	Therefore, $D_v$ is a linear map ($D_v$ is $\R$-linear). Furthermore, we know that $D_v$ satisfies the Leibniz rule
	\begin{gather*}
		D_v (fg) = (D_v f) g(p) + f(p) D_v g
	\end{gather*}
	since
	\begin{gather*}
		D_v (fg) = \sum_{i = 1}^n v^i \dfrac{\partial (fg)}{\partial x^i}(p) = \sum_{i = 1}^n v^i \left( \dfrac{\partial f}{\partial x^i}(p) g(p) + f(p) \dfrac{\partial g}{\partial x^i}(p) \right) \\
		= g(p) \sum_{i = 1}^n v^i \dfrac{\partial f}{\partial x^i}(p) + f(p) \sum_{i = 1}^n v^i \dfrac{\partial g}{\partial x^i}(p) = (D_v f) g(p) + f(p) (D_v g)
	\end{gather*}
	
	Any linear map $D: \Cinfty_p \to \R$ that satisfies the Leibniz rule is referred to as a \textit{derivation at $p$} or a \textit{point-derivation} of $\Cinfty_p$. Denote the set of all derivations at $p$ by $\mathcal{D}_p(\R^n)$. This space is a vector space.
	
	We know that all directional derivatives at $p$ are all derivations at $p$, so there exists a map
	\begin{gather*}
		\phi: T_p (\R^n) \to \mathcal{D}_p(\R^n)
	\end{gather*}
	given by
	\begin{gather*}
		v \mapsto D_v = \sum_{i = 1}^n v^i \dfrac{\partial}{\partial x^i} \Bigg|_p
	\end{gather*}
	We can deduce that $\phi$ is a linear map since for vectors $v, w \in T_p (\R^n)$ and $r \in \R$
	\begin{gather*}
		\phi(v + w) = D_{v + w} = \sum_{i = 1}^n (v + w)^i \dfrac{\partial}{\partial x^i} \Bigg|_p = \sum_{i = 1}^n (v^i + w^i) \dfrac{\partial}{\partial x^i} \Bigg|_p \\
		= \sum_{i = 1}^n v^i \dfrac{\partial}{\partial x^i} \Bigg|_p + \sum_{i = 1}^n w^i \dfrac{\partial}{\partial x^i} \Bigg|_p = D_v + D_w = \phi(v) + \phi(w)
	\end{gather*}
	and
	\begin{gather*}
		\phi(rw) = D_{rw} = \sum_{i = 1}^n (rw)^i \dfrac{\partial}{\partial x^i} \Bigg|_p = \sum_{i = 1}^n r(w^i) \dfrac{\partial}{\partial x^i} \Bigg|_p \\
		= r \sum_{i = 1}^n w^i \dfrac{\partial}{\partial x^i} \Bigg|_p = r D_w = r \phi(w)
	\end{gather*}
	
	\subsubsection*{Lemma 2.1}
	If $D$ is a point-derivation of $\Cinfty_p$, then $D(c) = 0$ for any constant function $c$.
	
	\noindent \textit{Proof}: Since every point-derivation is $\R$-linear, we know that $D(c) = cD(1)$, and so what we need to show is that $D(1) = 0$. We know by Leibniz rule that
	\begin{gather*}
		D(1) = D(1 \cdot 1) = D(1) \cdot 1 + 1 \cdot D(1) = 2D(1)
	\end{gather*}
	and so $D(1) = 0$. Thus, $D(c) = 0$.
	
	The Kronecker delta $\delta$ is a convenient notation used defined in the following manner
	\begin{gather*}
		\delta_i^j = \begin{cases}
			1 \quad i = j, \\
			0 \quad i \neq j
		\end{cases}
	\end{gather*}
	
	\subsubsection*{Theorem 2.2}
	The linear map $\phi: T_p (\R^n) \to \mathcal{D}_p(\R^n)$ given by $v \mapsto D_v$ is an isomorphism of vector spaces.
	
	\noindent \textit{Proof}: To show that this is an isomorphism, we have to show that it is bijective and that it is a linear map. We have already shown it is a linear operator, so now, we must show it is bijective.
	
	To show that it is injective, we need to show that $\phi(a) = \phi(b)$ implies $a = b$. By using the linear properties of $\phi$, we have that
	\begin{gather*}
		\phi(a) - \phi(b) = \phi(a - b) = 0
	\end{gather*}
	and so $a - b \in \text{ker}(\phi)$. To show that $a = b$, we need to show that the only vector in the kernel of $\phi$ is the zero vector. Suppose that for some vector $v \in T_p (\R^n)$, we have $D_v = 0$, i.e. $v$ is a vector in the kernel of $\phi$. Let $x^j : \R^n \to \R$ be the coordinate function sending a point $q$ to its coordinate $q^j$. Applying $D_v$ to the coordinate function, we have that
	\begin{gather*}
		D_v x^j = \sum_{i = 1}^n v^i \dfrac{\partial}{\partial x^i} \Bigg|_p x^j = \sum_{i = 1}^n v^i \delta_i^j = v^j = 0
	\end{gather*}
	Therefore, $v = 0$ and so the kernel of $\phi$ consists solely of the zero vector, meaning $\phi$ is injective.
	
	To show that it is surjective, let $D$ be a derivation at $p$, i.e. $D \in \mathcal{D}_p(\R^n)$ and let $(f, V)$ be a representative of a germ in $\Cinfty_p$. By making $V$ as small as necessary, we can assume that $V$ is an open ball. By Taylor's Theorem with remainder, we know there exists $\Cinfty$ functions $g_1(x), \cdots, g_n(x)$ in a neighborhood of $p$ such that
	\begin{gather*}
		f(x) = f(p) + \sum_{i = 1}^n (x^i - p^i) g_i(x), \quad g_i(p) = \dfrac{\partial f}{\partial x^i}(p)
	\end{gather*}
	Apply $D$ to both sides
	\begin{gather*}
		D f(x) = D \left( f(p) + \sum_{i = 1}^n (x^i - p^i) g_i(x) \right) = D f(p) + D \left( \sum_{i = 1}^n (x^i - p^i) g_i(x) \right) \\
		= D f(p) + \sum_{i = 1}^n D \left( (x^i - p^i) g_i(x) \right)
	\end{gather*}
	From Leibniz rule, we know that
	\begin{gather*}
		D \left( (x^i - p^i) g_i(x) \right) = D (x^i - p^i) g_i(p) + (p^i - p^i) D g_i(x) \\
		= (D x^i) g_i(p) - (D p^i) g_i(p)
	\end{gather*}
	Since $f(p)$ and $p^i$ are constants, $D f(p) = Dp^i = 0$ and so
	\begin{gather*}
		D f(x) = \sum_{i = 1}^n (D x^i) g_i(p) = \sum_{i = 1}^n (D x^i) \dfrac{\partial f}{\partial x^i}(p)
	\end{gather*}
	Thus, there exists a vector $v = \langle D x^1, \cdots, D x^n \rangle$ such that $D = D_v$. \textit{Intuitively, a derivation $D$ at a point $p$ is completely determined by how it acts on the coordinate functions $x^1, \cdots, x^n$. The values $D x^i$ give the components of the vector $v$, showing that $D$ is just the directional derivative in that direction.}

	This theorem allows us to identify the tangent vectors at $p$ with the derivations at $p$. Under the isomorphism $T_p (\R^n) \cong \mathcal{D}_p(\R^n)$, the standard basis $e_1, \cdots, e_n$ for $T_p (\R^n)$ corresponds to the set $\partial / \partial x^1 |_p, \cdots, \partial / \partial x^n |_p $ of partial derivatives. Thus, we can write a tangent vector $v = \langle v^1, \cdots, v^n \rangle = \sum v^i e_i$ as
	\begin{gather*}
		v = \sum v^i \dfrac{\partial}{\partial x^i} \Bigg|_p
	\end{gather*}
	
	While $T_p (\R^n)$ is isomorphic to $\mathcal{D}_p(\R^n)$, there is a reason the latter is preferred over the former, which is that a vector $v \in T_p (\R^n)$ depends on a choice of coordinates to describe it, while a derivation $D \in \mathcal{D}_p(\R^n)$ is independent of coordinates and one solely relies on how it impacts functions to determine its behavior.
	
	\subsection*{Vector Fields}
	A \textit{vector field} $X$ on an open subset $U \subset \R^n$ is a function that assigns to each point $p$ a tangent vector $X_p$ in $T_p (\R^n)$. Since $T_p (\R^n)$ is isomorphic to $\mathcal{D}_p(\R^n)$, it is convenient in this instance to use the basis $\{ \partial / \partial x^i |_p \}$. Thus the vector $X_p$ is a linear combination
	\begin{gather*}
		X_p = \sum a^i(p) \dfrac{\partial}{\partial x^i} \Bigg|_p
	\end{gather*}
	We can omit $p$, thus expressing $X$ as
	\begin{gather*}
		X = \sum a^i \dfrac{\partial}{\partial x^i}
	\end{gather*}
	where $a^i$ are functions on $U$. We say the vector field $X$ is $\Cinfty$ on $U$ if the coefficient functions $a^i$ are $\Cinfty$ on $U$.

	Consider the space $\R^2 - \{0\}$ and let $p = (x, y)$. Then we can define a vector field $X$ on this space by
	\begin{gather*}
		X = \dfrac{-y}{\sqrt{x^2 + y^2}} \dfrac{\partial }{\partial x} + \dfrac{x}{\sqrt{x^2 + y^2}} \dfrac{\partial }{\partial y} = \left\langle \dfrac{-y}{\sqrt{x^2 + y^2}}, \dfrac{x}{\sqrt{x^2 + y^2}} \right\rangle
	\end{gather*}
	The vector $X_p$ at a point $p$ is a unit vector on a circle of radius $\sqrt{x^2 + y^2}$ perpendicular to it going counter. 
	
	For another example, on the space $\R^2$, consider the vector field
	\begin{gather*}
		X = \langle x, -y \rangle
	\end{gather*}
	The vector field, when visualized by plotting various vectors $X_p$ attached to different points $p$, looks to move towards the $x$ axis and away from the $y$ axis.
	
	One can identify the vector fields on $U$ with column vectors of $\Cinfty$ functions on $U$, that is
	\begin{gather*}
		X = \sum a^i \dfrac{\partial}{\partial x^i} \iff
		\begin{bmatrix}
			a^1 \\
			\vdots \\
			a^n
		\end{bmatrix}
	\end{gather*}
	Notice here that it is very similar to the reinterpretation of the vector
	\begin{gather*}
		v = \sum v^i \dfrac{\partial}{\partial x^i} \Bigg|_p
	\end{gather*}
	with $v^i$ (a number corresponding to a point $p$) changing to $a^i$ (a function corresponding to $U$, becoming a number only when the specific point is chosen).
	
	The set of $\Cinfty$ functions equipped with the standard addition function obviously forms a ring, typically denoted by $\Cinfty(U)$ or $\mathcal{F}(U)$. We will demonstrate it is a ring in Exercise 2.2.
	
	\textit{Tu doesn't make it clear that he is no longer working with $\Cinfty(U, \R^m)$ functions, i.e. functions from $U \subset \R^n$ to $\R^m$, and has shifted to $\Cinfty(U, \R)$, i.e. functions from $U \subset \R^n$ to $\R$. This shift is important to keep in mind to understand the context. Importantly, without keepining this in mind, the only way to establish that $\Cinfty(U, \R^m)$ is a ring is to utilize component-wise multiplication, a very controversial choice that I personally found uncomfortable.}
	
	Multiplication of vector fields by functions on $U$ is defined pointwise
	\begin{gather*}
		(fX)_p = f(p) X_p, \quad p \in U
	\end{gather*}
	Since $X = \sum a^i \partial / \partial x^i$ is a $\Cinfty$ vector field and $f$ is a $\Cinfty$ function on $U$, $fX = \sum (fa^i) \partial / \partial x^i$ is also a vector field. Thus, the set of all $\Cinfty$ vector fields on $U$, denoted by $\mathbb{X}(U)$, is not just a vector space over $\R$ but also a module over the ring $\Cinfty(U)$ (but crucially, it is NOT a vector space over $\Cinfty(U)$ since $\Cinfty(U)$ is not a field).
	
	\subsubsection*{Module}
	If $R$ is a commutative ring (not rng), then a (left) $R$-module is an abelian group $A$ with scalar multiplication
	\begin{gather*}
		\mu : R \times A \to A
	\end{gather*}
	typically denoted $\mu(r, a) = ra$, such that for $r, s \in R$ and $a, b \in A$
	\begin{enumerate}
		\item (Associativity) $(rs)a = r(sa)$
		\item (Identity) If $1$ is the multiplicative identity in $R$, then $1a = a$
		\item (Distributivity) $(r + s)a = ra + sa$ and $r(a + b) = ra + rb$
	\end{enumerate}
	
	\subsubsection*{Module Homomorphism}
	Let $A$ and $A'$ be $R$-modules. An $R$-module homomorphism between $A$ and $A'$ is a map $f : A \to A'$ that preserves addition and scalar multiplication, i.e. for all $a, b \in A$ and $r \in R$
	\begin{enumerate}
		\item $f(a + b) = f(a) + f(b)$
		\item $f(ra) = rf(a)$
	\end{enumerate}
	
	\subsubsection*{Vector Fields as Derivations}
	If $X$ is a $\Cinfty$ vector field on an open subset $U \subset \R^n$ and $f$ is $\Cinfty$ function on $U$, then we define a new function $Xf$ on $U$ by
	\begin{gather*}
		(Xf)(p) = X_p f, \quad p \in U
	\end{gather*}
	Recall, though, that $X = \sum a^i \partial / \partial x^i$, meaning
	\begin{gather*}
		(Xf)(p) = \sum a^i(p) \dfrac{\partial f}{\partial x^i}(p)
	\end{gather*}
	and so
	\begin{gather*}
		Xf = \sum a^i \dfrac{\partial f}{\partial x^i}
	\end{gather*}
	a $\Cinfty$ function on $U$. We can thus interpret the vector field $X$ as an $\R$-linear map
	\begin{gather*}
		\Cinfty(U) \mapsto \Cinfty(U) \\
		f \mapsto Xf
	\end{gather*}
	And from here, we can show that $X$ satisfies the Leibniz rule, thus establishing that it is a derivation.	
	
	\subsubsection*{Leibniz rule for vector fields}
	Given that $X$ is a $\Cinfty$ vector field and $f, g$ are $\Cinfty$ functions on an open set $U \subset \R^n$, then $X(fg)$ satisfies the Leibniz rule
	\begin{gather*}
		X(fg) = (Xf)g + f(Xg)
	\end{gather*}
	
	\noindent \textit{Proof}: $X(fg)$ is defined on $U$ by
	\begin{gather*}
		(X(fg))(p) = X_p(fg)
	\end{gather*}
	for all $p \in U$. We know that
	\begin{gather*}
		X_p(fg) = \sum a^i(p) \dfrac{\partial (fg)}{\partial x^i} (p) = \sum a^i(p) \left( \dfrac{\partial f}{\partial x^i} (p) g(p) + f(p) \dfrac{\partial g}{\partial x^i} (p) \right) \\
		= \sum \left( a^i(p) \dfrac{\partial f}{\partial x^i} (p) \right) g(p) + \sum \left( a^i(p) \dfrac{\partial g}{\partial x^i} (p) \right) f(p) \\
		= \left( \sum a^i(p) \dfrac{\partial f}{\partial x^i} (p) \right) g(p) + \left( \sum a^i(p) \dfrac{\partial g}{\partial x^i} (p) \right) f(p) \\
		= (X_p f) g(p) + f(p) (X_p g)
	\end{gather*}
	and so 
	\begin{gather*}
		X(fg) = (Xf)g + f(Xg)
	\end{gather*}
	
	\subsubsection*{Generalization of the Concept of Derivation}
	If $A$ is an algebra over a field $K$, the derivation of $A$ is a $K$-linear map $D: A \to A$ such that $D$ satisfies the Leibniz rule
	\begin{gather*}
		D(ab) = (Da)b + a(Db)
	\end{gather*}
	for all $a, b \in A$. The set of all derivations of $A$, denoted by $\text{Der}(A)$, is closed under addition and scalar multiplication.
	
	We have just shown that a $\Cinfty$ vector field $X$ gives rise to a derivation of the algebra $\Cinfty(U)$ (and not just $\Cinfty_p$). Now consider the map
	\begin{gather*}
		\phi: \mathbb{X} \to \text{Der}(\Cinfty(U)) \\
		X \mapsto (f \mapsto (Xf))
	\end{gather*}
	Recall earlier that the initial conception of the vector field was a function that assigns to each point $p \in U$ a tangent vector $v \in T_p(\R^n)$. What this map informs us is that just like how we can identify tangent vectors $v$ at point $p$ with the point-derivations at $\Cinfty_p$ (as $T_p(\R^n)$ is isomorphic to $\mathcal{D}_p(\R^n)$), so to can vector fields on an open set $U$ be identified with the derivations of the algebra $\Cinfty(U)$ (we can treat it sort of like a \textit{flow} overlapping the space, how much \textit{work} is done by the vector field as an object traverses a path on $U$).
	
	\noindent \textit{Comment: Tu doesn't explicitly contrasts $\text{Der}(\Cinfty(U))$ with the set of all pointwise operators formed by arbitrarily gluing together point-derivations. However, we can construct a counterexample to show where they differ. Consider the space $\R$. At each point $p \in \R$, we select a point derivation $\tilde{D}(p)$ via the following assignment}
	\begin{gather*}
		\tilde{D}(p) = \begin{cases}
			\dfrac{\partial }{\partial x}, \quad p \in \mathbb{Q} \\
			-\dfrac{\partial }{\partial x}, \quad p \notin \mathbb{Q}
		\end{cases}
	\end{gather*}
	\textit{We are entirely free to glue together all these point-derivations to get a pointwise operator $\tilde{D} \in \prod_{p \in \mathbb{R}} \mathcal{D}_p(\mathbb{R})$. However, this operator is not an element in $\text{Der}(\Cinfty(U))$, because if it did, its component function $a(p)$ would have to be $\Cinfty$, but here we see that this is not the case.}
	
	
	\subsubsection*{Exercise 2.1}
	Let $X = x \partial / \partial x + y \partial / \partial y$ be a vector field on $\R^3$ and let $f(x, y, z) = x^2 + y^2 + z^2$ be a function on $\R^3$. Compute $Xf$
	
	Recall that $Xf$ is given by
	\begin{gather*}
		Xf = \sum a^i \dfrac{\partial f}{\partial x^i}
	\end{gather*}
	We see that $a^1 = x$, $a^2 = y$, and $a^3 = 0$, and so
	\begin{gather*}
		X_p f = x \dfrac{\partial f}{\partial x} \Bigg|_p + y \dfrac{\partial f}{\partial y} \Bigg|_p + 0 \dfrac{\partial f}{\partial z} \Bigg|_p 
		= x \cdot 2 x + y \cdot 2 y = 2 \left( x^2 + y^2 \right)
	\end{gather*}
	
	\subsubsection*{Exercise 2.2}
	Define addition, multiplication, and scalar multiplication in $\Cinfty_p$. Prove that addition in $\Cinfty_p$ is commutative.
	
	Recall that $\Cinfty_p$ is the set of all germs of $\Cinfty$ functions on $\R^n$ at $p$. Furthermore, recall that two pairs $(f, U)$ and $(g, V)$ (where $U$ and $V$ are neighborhoods of $p$) are equivalent to each other if there exists an open set $W \subset U \cap V$ containing $p$ such that $f|_W = g|_W$, and that the equivalence class of such pairs is called the germ of $f$ at $p$.
	
	Consider two germs $[f], [g] \in \Cinfty_p$. We know that there exists pairs
	\begin{gather*}
		(f_1, U_1) \in [f] \\
		(g_1, V_1) \in [g]
	\end{gather*}
	and since $U_1$ and $V_1$ are neighborhoods of $p$, we know there exists a nonempty set $W_1 \subset U_1 \cap V_1$ containing $p$. We define $[f] + [g]$ to be the equivalence class of all pairs equivalent to $(f_1|_{W_1} + g_1|_{W_1}, W_1)$. In a similar manner, we define $[f] \cdot [g]$ to be the equivalence class of all pairs equivalent to  $(f_1|_{W_1} \cdot g_1|_{W_1}, W_1)$
	
	To show that addition and multiplication is well-defined, consider two other pairs
	\begin{gather*}
		(f_2, U_2) \in [f] \\
		(g_2, V_2) \in [g]
	\end{gather*}
	Since $U_2$ and $V_2$ are neighborhoods of $p$, we know there exists a nonempty set $W_2 \subset U_2 \cap V_2$ containing $p$. Finally, since both $W_1$ and $W_2$ are neighborhoods of $p$, there must exist a nonempty set $W \subset W_1 \cap W_2$ containing $p$. Since both $(f_1, U_1)$ and $(f_2, U_2)$ are elements of the same germ of $f$ at $p$, we know that $f_1|_U = f_2|_U$ for some set $U$. Similarly, since both $(g_1, V_1)$ and $(g_2, V_2)$ are elements of $[g]$, we know that $g_1|_V = g_2|_V$ for some set $V$. Now, consider the set $W = W_1 \cap W_2 \cap U \cap V$. Clearly, $W \subset U_1 \cap U_2$ and $W \subset V_1 \cap V_2$, and so
	\begin{gather*}
		f_1|_W = f_2|_W \\
		g_1|_W = g_2|_W
	\end{gather*}
	and from here, we see that for addition, we have
	\begin{gather*}
		\left( f_1|_{W_1} + g_1|_{W_1} \right)|_W = f_1|_W + g_1|_W = f_2|_W + g_2|_W
		= \left( f_2|_{W_2} + g_2|_{W_2} \right)|_W
	\end{gather*}
	and for multiplication, we have
	\begin{gather*}
		\left( f_1|_{W_1} \cdot g_1|_{W_1} \right)|_W = f_1|_W \cdot g_1|_W
		= f_2|_W \cdot g_2|_W = \left( f_2|_{W_2} \cdot g_2|_{W_2} \right)|_W	
	\end{gather*}
	Therefore $(f_1|_{W_1} + g_1|_{W_1}, W_1)$ and $(f_2|_{W_2} + g_2|_{W_2}, W_2)$ are part of the same equivalence class $[f] + [g]$, and $(f_1|_{W_1} \cdot g_1|_{W_1}, W_1)$ and $(f_2|_{W_2} \cdot g_2|_{W_2}, W_2)$ are part of the same equivalence class $[f] \cdot [g]$. Addition and multiplication are commutative since
	\begin{gather*}
		\left( f_1|_{W_1} + g_1|_{W_1} \right)(x) = f_1|_{W_1}(x) + g_1|_{W_1}(x) \\
		= g_1|_{W_1}(x) + f_1|_{W_1}(x) = \left( g_1|_{W_1} + f_1|_{W_1} \right)(x) 
	\end{gather*}
	and
	\begin{gather*}
		\left( f_1|_{W_1} \cdot g_1|_{W_1} \right)(x) = f_1|_{W_1}(x) \cdot g_1|_{W_1}(x) \\
		= g_1|_{W_1}(x) \cdot f_1|_{W_1}(x) = \left( g_1|_{W_1} \cdot f_1|_{W_1} \right)(x) 
	\end{gather*}
	
	Finally, let $c \in \R$, $[f] \in \Cinfty_p$, and $(f, U) \in [f]$. Then $c [f]$ is defined to be the equivalence class of all pairs equivalent to $(c f, U)$. Scalar multiplication is well-defined because for pairs $(f_1, U_1), (f_2, U_2) \in [f]$, there exists a set $U \subset U_1 \cap U_2$ such that $f_1|_U = f_2|_U$ and so $cf_1|_U = cf_2|_U$. Therefore, $(cf_1, U_1)$ and $(cf_2, U_2)$ are part of the same equivalence class.
	
\end{document}